\section{Limits}\label{sec:limits}

This section fixes what cannot be done, so that the positive results of
\cref{sec:envelopes,sec:part1,sec:part2} can be read against it.

\subsection{Without a prover: structure is necessary}

\begin{proposition}[worst case]\label{prop:worstcase}
\begin{enumerate}[label=(\alph*)]
\item If a sound efficient procedure is complete (for any fixed $\eps<1$) on
all 3-CNF formulas under the uniform distribution, then $\NP\subseteq\BPP$.
\item Let $G_k:\bits^k\to\bits^{2k}$ be computable by $\poly(k)$-size
circuits and suppose every circuit of size $T(k)$ distinguishes $G_k(U_k)$
from $U_{2k}$ with advantage less than $1/2$. For $y\in\bits^{2k}$ let
$C_y(x)=[G_k(x)=y]$ on inputs $x\in\bits^k$. Then $p(C_{G(s)})\ge 2^{-k}$ for
every seed $s$, while $\Pr_y[p(C_y)>0]\le2^{-k}$ for uniform $y$; and any
sound procedure with $\delta\le1/8$ that is complete on $\{C_{G(s)}\}$ runs in
time at least $T(k)-\poly(k)$ on some $C_{G(s)}$.
\end{enumerate}
\end{proposition}

\begin{proof}
(a) If $p=0$, a tight interval must be $[0,0]$; if $p>0$ (so $p\ge2^{-n}$) a
tight interval that contains $p$ has $L>0$. Answering ``satisfiable'' iff $L>0$
errs with probability at most $\delta$.
(b) The bounds on $p$ are immediate ($\E_y|G^{-1}(y)|=2^{k-2k}$). Let $A$ be
the procedure and $t$ its running time on the family. Consider the test
``$A(C_y)$ outputs $L>0$''. On $y=G(s)$ it fires with probability at least
$1-\delta$ (a tight interval around $p>0$ has $L>0$). On uniform $y$, with
probability at least $1-2^{-k}$ we have $p=0$ and soundness forces $L\le
p=0$ except with probability $\delta$. The advantage is at least
$1-2\delta-2^{-k}>1/2$; fixing $A$'s coins to their best value gives a circuit
of size $t+\poly(k)$, so $t\ge T(k)-\poly(k)$.
\end{proof}

Monte Carlo on $C_{G(s)}$ costs $O(2^k\poly(k))$. So if $G$ is
exponentially hard ($T(k)=2^{(1-o(1))k}$), no sound procedure beats sampling on
this family by more than a $2^{o(k)}=p^{-o(1)}$ factor, even though every
instance is a small, explicit circuit. Beating sampling is a matter of
structure: the classes of \cref{sec:part1} are where the structure is visible to
a polynomial-time algorithm, and those of \cref{sec:part2} are where it is
visible to a polynomial-time verifier given a hint.

\subsection{With a prover: upper bounds need deduction}

Lower and upper bounds behave very differently. A lower bound $p\ge L$ can be
witnessed by \emph{showing} probability mass: a proposal distribution $Q$
concentrated on $C^{-1}(1)$ (\cref{sec:lower}). An upper bound must exclude
mass that the verifier has not seen. The next theorem makes this precise:
a verifier that only evaluates $C$ cannot certify an upper bound much below
the reciprocal of its number of evaluations, whatever argument it is given.

\begin{theorem}[sampling cannot certify upper bounds]\label{thm:blackbox}
Let $V$ be a verifier that receives an argument $\pi$, has oracle access to
$f:\bits^n\to\bits$, flips coins, makes at most $T$ oracle queries (at
adaptively chosen points, from any distributions it likes), and outputs a
number $U$. Suppose $V$ is sound for upper bounds under the uniform
distribution: for all $f,\pi$, $\Pr[U<p(f)]\le\delta$. Then for every $f$,
every $\pi$ and every $k\le n$,
\[
\Pr\big[V^f(\pi)\text{ outputs }U<2^{-k}\big]\;\le\;\delta+T\,2^{-k}.
\]
In particular, to output an upper bound below $\eta$ with probability
$1-\delta'$ a black-box verifier needs $T\ge(1-\delta-\delta')/(2\eta)$
queries.
\end{theorem}

\begin{proof}
Fix $f$, $\pi$, $k$. For $t\in\bits^k$ let $H_t=\{x:x_{1..k}=t\}$ and
$f_t=f\vee\one_{H_t}$, so $p(f_t)\ge2^{-k}$. Fix the coins $r$ of $V$. The run
of $V^f(\pi;r)$ queries a fixed sequence of points $z_1,\dots,z_T$. The run of
$V^{f_t}(\pi;r)$ is identical unless some $z_i\in H_t$; for each fixed $z_i$
this happens for exactly one $t$. Hence for uniform $t$,
$\Pr_{r,t}[\text{the runs differ}]\le T2^{-k}$, and
\[
\E_t\Pr_r\big[V^{f_t}(\pi;r)<2^{-k}\big]\;\ge\;\Pr_r\big[V^{f}(\pi;r)<2^{-k}\big]-T2^{-k}.
\]
Some $t$ attains at least the average; for it, $V$ outputs $U<2^{-k}\le
p(f_t)$ with that probability, which soundness bounds by $\delta$.
For the last claim take $k=\lceil\log_2(1/\eta)\rceil-1$, so that
$\eta\le 2^{-k}<2\eta$.
\end{proof}

\begin{remark}
The same argument works for any distribution $D$ with enough min-entropy,
planting $H=\{x:Ax=b\}$ for a random $k\times n$ matrix $A$ over $\F_2$ and
random $b$: every fixed point lies in $H$ with probability $2^{-k}$, and by
pairwise independence $D(H)\ge2^{-k-1}$ with probability $1-4\cdot2^{k}\sum_x
D(x)^2$. The planted circuit $f\vee[Ax=b]$ is only $O(nk)$ gates larger, so
the theorem also holds for white-box verifiers that look at $C$ only through
evaluations and its size.
\end{remark}

The verifiers of \cref{sec:part2} escape \cref{thm:blackbox} by reading the
circuit: they check a proof that $C(x)=1$ forces $x$ into a region of small,
exactly computable measure. In general such proofs cannot always exist:

\begin{theorem}[no complete upper-bound arguments]\label{thm:ma}
Suppose a sound argument system (soundness error $\delta\le1/3$) is complete
for upper bounds on all 3-CNF formulas with $p=0$ under the uniform
distribution, i.e.\ for each such $C$ some $\pi$ of polynomial size makes $V$
output $U<2^{-n}$ with probability at least $2/3$. Then $\coNP\subseteq\MA$,
and the polynomial hierarchy collapses to its second level.
\end{theorem}

\begin{proof}
Merlin sends $\pi$; Arthur runs $V$ and accepts iff $U<2^{-n}$. If $C$ is
unsatisfiable he accepts with probability $\ge2/3$. If $C$ is satisfiable then
$p\ge2^{-n}$, and accepting means $V$ output $U<p$, which happens with
probability at most $1/3$. So 3-UNSAT$\in\MA\subseteq\AM$, and
$\coNP\subseteq\AM$ collapses $\PH$ to $\Sigma_2^p$ \cite{bhz1987}.
\end{proof}

The lower-bound side is weaker but not free: a system complete for lower
bounds within a constant factor on all circuits is exactly a proof that the
class $\SBP$ of \cite{bgm2006} is contained in $\MA$; since
$\MA\subseteq\SBP\subseteq\AM$ and the inclusions are not known to collapse,
we should not expect proposal-based lower bounds for every circuit either.

\subsection{Interaction trivialises the problem}

The restriction to one message is essential:

\begin{proposition}\label{prop:interactive}
\begin{enumerate}[label=(\alph*)]
\item (Sum-check \cite{lfkn1992}.) For every circuit $C$ and product $D$ with
dyadic marginals, the exact value of $p$ has an interactive proof with a
polynomial-time verifier, $O(n+|C|)$ rounds and soundness error
$O((n+|C|)\,d/|\F|)$, where $d$ bounds the individual degrees of the
arithmetised Tseitin constraint polynomial.
\item (Set lower bounds \cite{gs1986}.) Two-round public-coin protocols certify
$p\ge L$ within a constant factor for every circuit.
\end{enumerate}
\end{proposition}

Applying the Fiat--Shamir transform to (a) gives a one-message argument that
is only computationally sound, which is not the model asked for. So the
interesting question is the Merlin--Arthur one, where by
\cref{thm:blackbox,thm:ma} an argument must contain \emph{deductive} content
about the circuit, and can be complete only on a restricted class.
